\section{Discussion}

Here, we discuss the physical and architectural parameters that affect the performance of our system, highlight its advantages, and point out its limitations.

\subsection{Analysis of System Performance and Design Decisions}

Our experiments show that the proposed system achieves high accuracy, good transferability, and high computational efficiency. We can attribute these results to three key decisions: the expanded descriptor set, bidirectional temporal modeling, and attention based sequence aggregation.

\subsubsection{Kinematic Representational Capacity}
The primary factor driving the improvement in classification accuracy is the representational capacity of the nine-feature optical-flow descriptor. The feature ablation study (Section~5.2) showed a 5.24\% improvement in F1-score when moving from the three baseline features (Entropy, TOV, KDE) to the proposed nine-feature set. 

Physically, crowd panic and the onset of a crowd crush are characterized by a sudden loss of directional coordination and a wide dispersion of velocities as individuals attempt to escape in different directions. While directional entropy captures the overall disorder of movement directions, it is insensitive to the degree of alignment among adjacent motion vectors. The inclusion of Direction Variance ($\sigma^2_\theta$) and Motion Coherence ($C$) addresses this limitation by providing explicit measures of angular agreement. Furthermore, the inclusion of higher-order moments of the velocity distribution---specifically, the Standard Deviation ($\sigma_M$) and Skewness ($\gamma_M$) of the flow magnitudes---enables the model to distinguish between a crowd moving rapidly in a coordinated manner (e.g., normal egress where velocity variance is low and skewness is stable) and the early stages of panic (where velocity variance increases as some individuals run while others freeze).

\subsubsection{Bidirectional Temporal Context}
The comparison of temporal classifiers (Section~5.1) confirmed that the Bidirectional LSTM outperforms its unidirectional counterpart. Unidirectional recurrent networks must classify the current sequence using only historical states ($\mathbf{h}_{1:t}$). However, during the initiation phase of a stampede, the motion descriptors of the first few transitional frames can be statistically similar to transient normal variations (e.g., a localized group of people rushing to catch a train). 

By processing the temporal sequence in both forward and reverse directions, the BiLSTM generates hidden representations that incorporate context from both preceding and succeeding states. Consequently, a temporary velocity spike that is followed by a return to normal flow is represented differently from a velocity spike that is followed by escalating multidirectional motion. This bidirectional context reduces false positive alerts, contributing to the 100\% precision achieved by the BiLSTM architectures.

\subsubsection{Saliency-Guided Aggregation}
The additive attention layer helps to address the dilution of temporal signals in standard recurrent networks. In typical crowd safety systems, anomalous events occupy a fraction of the temporal window, and the rest of the window contains normal patterns. Without an attention mechanism, classification is done using a fixed vector, which dilutes the signatures of transient anomalies.

The attention mechanism assigns non uniform weights to every time step depending on its importance for classification. The attention maps show that the attention weights concentrate exactly on the transition windows where the motion features begin to deviate from normal values. Such selection improves the classification accuracy and provides temporal localization of the event.

\subsection{Operational Strengths of the Proposed Pipeline}

The proposed system exhibits four characteristics that are relevant for practical deployment:

\begin{enumerate}
  \item \textbf{Low Computational Footprint:} The total model size is 554,817 parameters and processing takes 18.35 ms per frame, which is equal to 54.5 FPS throughput on a standard laptop GPU. This is much smaller than deep spatial architectures, making it possible to run the system on resource-constrained edge computing devices near the cameras, without the need for dedicated servers.
  \item \textbf{Explainability:} The permutation feature importance and attention weight visualization give post-hoc explanations at the feature level and temporal level. This interpretability is crucial for safety-critical environments, where operators must verify automatic alerts before responding.
  \item \textbf{Cross-Environment Generalisation:} The zero-shot testing on the UCSD Ped2 and CUHK Avenue datasets showed that the model generalizes to new environments without retraining. This is achieved by using scene-independent optical-flow statistics, which are robust to camera perspective, lighting, and layout changes.
  \item \textbf{Early Warning Capability:} The system detected anomalies 0.39~s before the annotated peak intensity on average, with the warning lead time reaching up to 7.5~s. In the crowd safety domain, even sub-second warnings are important since they can trigger automatic countermeasures, such as opening emergency doors, before the density of the crowd becomes dangerous.
\end{enumerate}

\subsection{Boundary Conditions and Limitations}

The system's performance is subject to several boundary conditions and limitations:

\begin{enumerate}
  \item \textbf{Sensitivity to Extreme Image Degradation:} As shown in Experiment 5, severe blur or noise causes a drop in the F1-Score below 55 percent. Since the feature extraction pipeline relies on spatial intensity gradients to estimate the optical flow, extreme degradation distorts the computed displacement vectors. The system relies on a minimum level of image quality and lighting to get reliable motion gradients.
  \item \textbf{Stationary Camera Assumption:} The dense optical flow computation assumes a stationary camera perspective. Camera movements, such as panning, tilting, or zooming, would introduce global displacement vectors across the entire frame, distorting the crowd motion statistics. The pipeline is designed for static surveillance cameras and would require global motion compensation techniques to operate on dynamic cameras.
  \item \textbf{Evaluation on Controlled Benchmarks:} While cross-dataset transfer was demonstrated, the evaluation was conducted using standard, controlled benchmark datasets. Actual surveillance installations often exhibit variable frame rates, high video compression ratios, and intermittent network latency. The performance of the system under these network-induced variations has not been evaluated.
  \item \textbf{Annotation Granularity:} The training data (UMN) lacks frame-level anomaly labels, requiring the use of a temporal heuristic (designating the final 30\% of each sequence as anomalous) to generate ground-truth labels. While this heuristic matches the overall structure of the videos, it introduces a degree of label noise at the boundary transition frames, which may introduce minor biases into the absolute metric values.
\end{enumerate}
